\section{近场边缘与孔径的空间频率}
夫琅禾费近似把孔面场变为空间频率的傅里叶变换；菲涅耳近似则保留每个孔元的二次相位。二者使用同一传播积分，区别在于可省略哪些相位。描述近场时，不能把每条射线的几何阴影边界当成强度突然为零的位置。

直边衍射需要的振荡积分也可由第一讲的复高斯积分确定。先乘收敛因子$\ee^{-\epsilon t^2}$，$\epsilon>0$，在式\eqref{eq:o1gaussintegral}中取$a=\epsilon-i\pi/2$、$b=0$，
\[
 \int_{-\infty}^{\infty}\ee^{-\epsilon t^2+i\pi t^2/2}\dd t
 =\sqrt{\frac{\pi}{\epsilon-i\pi/2}}
 \xrightarrow{\epsilon\to0^+}\sqrt2\ee^{i\pi/4}=1+i.
\]
平方根沿$\operatorname{Re}a>0$连续取支。该收敛极限给出菲涅耳振荡积分；被积函数为偶函数，所以$\int_0^\infty\ee^{i\pi t^2/2}\dd t=(1+i)/2$。实、虚部分别给$C(\infty)=S(\infty)=1/2$，用于下面的无遮挡场归一化。

\begin{example}\label{ex:o3edge}
平面波照射遮住$x<0$的直边，光沿$+z$传播。距边$L$的屏上坐标为$x$，定义$v=x\sqrt{2/(\lambda L)}$。用菲涅耳积分表示归一化强度，求几何边界$x=0$的强度。
\Diagram{fig:o3edgeq}{半平面边缘及几何阴影线}{\draw[curve](0,-1.5)--(0,0);\foreach\y in{-1.4,-1.1,-.8,-.5,-.2}{\draw(0,\y)--(-.3,{\y-.2});}\draw[curve](5,-1.8)--(5,2);\draw[guide](0,0)--(5,0);\draw[->](-1,1)--(1,1);\node[right]at(5,-1){几何阴影};\node[right]at(5,1){照明区};}
\end{example}
\begin{solution}
图\ref{fig:o3edgeq}的透光半平面为$\xi\geq0$。对沿边方向的积分提出共同因子后，垂直边方向场正比$\int_0^\infty\exp[i\pi(x-\xi)^2/(\lambda L)]\dd\xi$。令$t=(x-\xi)\sqrt{2/(\lambda L)}$，$\xi=0$对应$t=v$，$\xi\to\infty$对应$t\to-\infty$，负的微分将上下限反转，故积分正比$\int_{-\infty}^v\ee^{i\pi t^2/2}\dd t$。
\[
 C(v)=\int_0^v\cos(\pi t^2/2)\dd t,\quad
 S(v)=\int_0^v\sin(\pi t^2/2)\dd t.
\]
用$C(\infty)=S(\infty)=1/2$和奇函数性质，分子为$(1/2+C)+i(1/2+S)$。
\Diagram{fig:o3edgea}{按菲涅耳积分计算的直边衍射强度，几何边界处为四分之一}{\draw[axis](-3.2,0)--(3.7,0)node[right]{$v$};\draw[axis](0,0)--(0,3.2)node[above]{$I/I_0$};\draw[curve]plot coordinates {(-3.00000,0.012310) (-2.96425,0.012605) (-2.92849,0.012910) (-2.89274,0.013227) (-2.85698,0.013556) (-2.82123,0.013897) (-2.78547,0.014250) (-2.74972,0.014618) (-2.71397,0.014999) (-2.67821,0.015395) (-2.64246,0.015807) (-2.60670,0.016235) (-2.57095,0.016681) (-2.53520,0.017144) (-2.49944,0.017627) (-2.46369,0.018131) (-2.42793,0.018655) (-2.39218,0.019203) (-2.35642,0.019774) (-2.32067,0.020370) (-2.28492,0.020993) (-2.24916,0.021645) (-2.21341,0.022326) (-2.17765,0.023039) (-2.14190,0.023785) (-2.10615,0.024568) (-2.07039,0.025388) (-2.03464,0.026248) (-1.99888,0.027151) (-1.96313,0.028100) (-1.92737,0.029097) (-1.89162,0.030146) (-1.85587,0.031250) (-1.82011,0.032412) (-1.78436,0.033638) (-1.74860,0.034930) (-1.71285,0.036294) (-1.67709,0.037735) (-1.64134,0.039258) (-1.60559,0.040869) (-1.56983,0.042575) (-1.53408,0.044382) (-1.49832,0.046297) (-1.46257,0.048330) (-1.42682,0.050488) (-1.39106,0.052782) (-1.35531,0.055221) (-1.31955,0.057817) (-1.28380,0.060583) (-1.24804,0.063531) (-1.21229,0.066676) (-1.17654,0.070034) (-1.14078,0.073622) (-1.10503,0.077458) (-1.06927,0.081564) (-1.03352,0.085960) (-0.99777,0.090672) (-0.96201,0.095725) (-0.92626,0.101147) (-0.89050,0.106971) (-0.85475,0.113229) (-0.81899,0.119958) (-0.78324,0.127198) (-0.74749,0.134992) (-0.71173,0.143388) (-0.67598,0.152435) (-0.64022,0.162189) (-0.60447,0.172709) (-0.56872,0.184060) (-0.53296,0.196312) (-0.49721,0.209538) (-0.46145,0.223821) (-0.42570,0.239246) (-0.38994,0.255906) (-0.35419,0.273900) (-0.31844,0.293335) (-0.28268,0.314321) (-0.24693,0.336979) (-0.21117,0.361434) (-0.17542,0.387818) (-0.13966,0.416269) (-0.10391,0.446932) (-0.06816,0.479956) (-0.03240,0.515493) (0.00335,0.553700) (0.03911,0.594733) (0.07486,0.638752) (0.11061,0.685910) (0.14637,0.736357) (0.18212,0.790236) (0.21788,0.847678) (0.25363,0.908798) (0.28939,0.973691) (0.32514,1.042429) (0.36089,1.115053) (0.39665,1.191567) (0.43240,1.271932) (0.46816,1.356059) (0.50391,1.443804) (0.53966,1.534957) (0.57542,1.629234) (0.61117,1.726277) (0.64693,1.825637) (0.68268,1.926775) (0.71844,2.029055) (0.75419,2.131740) (0.78994,2.233992) (0.82570,2.334871) (0.86145,2.433339) (0.89721,2.528269) (0.93296,2.618454) (0.96872,2.702622) (1.00447,2.779462) (1.04022,2.847639) (1.07598,2.905834) (1.11173,2.952776) (1.14749,2.987280) (1.18324,3.008296) (1.21899,3.014955) (1.25475,3.006620) (1.29050,2.982935) (1.32626,2.943878) (1.36201,2.889802) (1.39777,2.821479) (1.43352,2.740125) (1.46927,2.647416) (1.50503,2.545493) (1.54078,2.436936) (1.57654,2.324730) (1.61229,2.212195) (1.64804,2.102900) (1.68380,2.000547) (1.71955,1.908834) (1.75531,1.831298) (1.79106,1.771145) (1.82682,1.731064) (1.86257,1.713059) (1.89832,1.718273) (1.93408,1.746857) (1.96983,1.797854) (2.00559,1.869154) (2.04134,1.957496) (2.07709,2.058541) (2.11285,2.167031) (2.14860,2.277011) (2.18436,2.382138) (2.22011,2.476043) (2.25587,2.552741) (2.29162,2.607058) (2.32737,2.635058) (2.36313,2.634413) (2.39888,2.604699) (2.43464,2.547578) (2.47039,2.466829) (2.50615,2.368210) (2.54190,2.259137) (2.57765,2.148196) (2.61341,2.044493) (2.64916,1.956893) (2.68492,1.893216) (2.72067,1.859441) (2.75642,1.859019) (2.79218,1.892370) (2.82793,1.956642) (2.86369,2.045788) (2.89944,2.150984) (2.93520,2.261397) (2.97095,2.365235) (3.00670,2.451002) (3.04246,2.508830) (3.07821,2.531726) (3.11397,2.516577) (3.14972,2.464759) (3.18547,2.382227) (3.22123,2.279016) (3.25698,2.168165) (3.29274,2.064152) (3.32849,1.981003) (3.36425,1.930323) (3.40000,1.919516)};\fill(0,.55)circle(1.5pt)node[left]{$1/4$};\draw[guide](-3,2.2)--(3.4,2.2);\node[right]at(3.4,2.2){$1$};}
无遮挡平面对应积分$-\infty<t<\infty$，值为$1+i$。因此图\ref{fig:o3edgea}使用的准确归一化为
\[
 \frac{U}{U_0}=\frac{(1/2+C)+i(1/2+S)}{1+i},\quad
 \frac I{I_0}=\frac{(1/2+C)^2+(1/2+S)^2}{2}.
\]
在$v=0$，分子为$(1+i)/2$，所以场为一半、强度为四分之一。$v\to-\infty$时$C,S\to-1/2$，深阴影强度趋零；$v\to+\infty$时均趋$1/2$，强度趋1。边界附近的空间尺度由$|v|\sim1$给出，即$|x|\sim\sqrt{\lambda L/2}$，所以屏越远，几何边缘被衍射展开的范围越大。正照、直边及均匀平面波是本归一化的条件。
\end{solution}

\begin{example}\label{ex:o3shiftproperty}
孔面场$U_0(\xi)$整体横移$\xi_0$。求夫琅禾费场的变化；若孔未移动而给孔面加线性相位$\exp(iq_0\xi)$，又怎样改变远场？
\Diagram{fig:o3shiftq}{孔的位置平移与孔内相位斜坡是两种不同操作}{\draw[axis](-1,0)--(5,0)node[right]{$\xi$};\draw[curve](0,0)rectangle(1.2,1);\draw[guide](2.5,0)rectangle(3.7,1);\draw[->](.6,1.5)--(3.1,1.5)node[midway,above]{$\xi_0$};}
\end{example}
\begin{solution}
取远场变量$q=k\sin\theta$，$F(q)=\int U_0(\xi)\ee^{-iq\xi}\dd\xi$。图\ref{fig:o3shiftq}平移孔为$U_0(\xi-\xi_0)$，令$s=\xi-\xi_0$，得$F_{\rm shift}(q)=\ee^{-iq\xi_0}F(q)$。单孔强度分布不变，只改变方向相关相位。
\Diagram{fig:o3shifta}{线性孔面相位使远场分布平移，而整体孔位置只加远场相位}{\draw[axis](-2.3,0)--(3.4,0)node[right]{$q$};\draw[curve,domain=-2:3,samples=70]plot(\x,{2*exp(-2*\x*\x)});\draw[guide,domain=-2:3,samples=70]plot(\x,{2*exp(-2*(\x-1.3)^2)});\node[above]at(1.3,2.1){$q_0$};}
若孔内乘$\ee^{iq_0\xi}$，积分变为$F(q-q_0)$，强度图样向$q_0$平移，如图\ref{fig:o3shifta}。斜入射照明的切向波矢或相位片的线性相位均按此式改变远场中心。孔整体平移在有另一参考孔相干叠加时会改变条纹，不能因单孔强度不变就删去其相位。
\end{solution}

\begin{example}\label{ex:o3twowidths}
两条中心相隔$b$的缝宽分别为$a_1,a_2$，单位面积照明场相同且无初始相差。求远场强度，说明可见度为何随角度变化。
\Diagram{fig:o3twowidthq}{不等缝宽带来不同角度包络和轴上振幅}{\draw[curve](0,-.6)rectangle(.6,1);\draw[curve](2,-.6)rectangle(3.2,1);\node[below]at(.3,-.6){$a_1$};\node[below]at(2.6,-.6){$a_2$};\draw[<->](.3,1.4)--(2.6,1.4)node[midway,above]{$b$};}
\end{example}
\begin{solution}
图\ref{fig:o3twowidthq}中每缝中心对称积分给实包络$F_j=a_j\operatorname{sinc}(\pi a_j\sin\theta/\lambda)$。合场$C[F_1+F_2\ee^{-ikb\sin\theta}]$，故$I=|C|^2[F_1^2+F_2^2+2F_1F_2\cos(kb\sin\theta)]$。
\Diagram{fig:o3twowidtha}{取$a_2=2a_1$时的有符号场包络，$\beta=\pi a_1\sin\theta/\lambda$}{
\draw[axis](-3.2,0)--(3.4,0)node[right]{$\beta$};
\draw[axis](0,-.65)--(0,2.7)node[above]{$F_j/a_1$};
\draw[curve,domain=-6.7:6.7,samples=181]plot({\x/2.3},{abs(\x)<.001 ? 1 : sin(deg(\x))/\x});
\draw[guide,domain=-6.7:6.7,samples=181]plot({\x/2.3},{abs(\x)<.001 ? 2 : sin(deg(2*\x))/\x});
\node[left]at(0,1){$1$};\node[left]at(0,2){$2$};
\node[below]at(1.366,0){$\pi$};\node[below]at(2.732,0){$2\pi$};
\draw[curve](1.6,2.2)--(2.2,2.2)node[right]{$F_1$};
\draw[guide](1.6,1.65)--(2.2,1.65)node[right]{$F_2$};}
在包络随一个纹距缓变的局部区域，$\mathcal V=2|F_1F_2|/(F_1^2+F_2^2)$。图\ref{fig:o3twowidtha}中两包络的零点不同：某一包络过零而另一不为零时，可见度降至0；过零后的$F_j$反号又使干涉纹相位翻转。若两者同时为零，该方向无光，不定义可见度。宽度比只给出轴上振幅比，离轴计算须保留各自的sinc因子。
\end{solution}

\section{连续孔径、二维光栅与采样}
\begin{example}\label{ex:o3circular}
半径$a$的均匀圆孔正入射。将远场场写成一维贝塞尔积分并求归一化形式，说明第一暗环如何决定角分辨尺度。
\Diagram{fig:o3circleq}{圆孔的径向对称性使远场只依赖离轴角}{\draw[curve](0,0)circle(1.6);\draw[axis](0,0)--(2.6,0)node[right]{$\xi$};\draw[axis](0,0)--(0,2.3)node[above]{$\eta$};\draw[->](0,0)--(1.1,1.1)node[right]{$r$};}
\end{example}
\begin{solution}
图\ref{fig:o3circleq}中，取观察方向的横向投影为极角零方向，$q=k\sin\theta$。孔面点的远场相位为$-qr\cos\varphi$，故
\[
 U=C\int_0^a r\dd r\int_0^{2\pi}\ee^{-iqr\cos\varphi}\dd\varphi
 =2\pi C\int_0^a rJ_0(qr)\dd r.
\]
这里用角积分定义$J_0(z)=(2\pi)^{-1}\int_0^{2\pi}\ee^{-iz\cos\varphi}\dd\varphi$。为求径向原函数，将指数展开；奇次余弦积分为零，偶次积分记为$C_{2j}=\int_0^{2\pi}\cos^{2j}\varphi\dd\varphi$。取$u=\cos^{2j-1}\varphi$、$\dd v=\cos\varphi\dd\varphi$，分部积分为
\begin{align*}
 C_{2j}&=[\cos^{2j-1}\varphi\sin\varphi]_0^{2\pi}
 +(2j-1)\int_0^{2\pi}\cos^{2j-2}\varphi\sin^2\varphi\dd\varphi\\
 &=(2j-1)(C_{2j-2}-C_{2j}),\qquad j\geq1.
\end{align*}
移项得$C_{2j}=(2j-1)C_{2j-2}/(2j)$，从$C_0=2\pi$递推，得$C_{2j}=2\pi(2j)!/[2^{2j}(j!)^2]$。于是
\[
 J_0(z)=\sum_{j=0}^{\infty}\frac{(-1)^j(z/2)^{2j}}{(j!)^2},
 \qquad
 J_1(z)=\sum_{j=0}^{\infty}
 \frac{(-1)^j(z/2)^{2j+1}}{j!(j+1)!}.
\]
第二式定义此处所需的一阶贝塞尔函数。逐项求导时，$zJ_1(z)$的第$j$项乘出$2j+2$，约去$(j+1)!$，便得到$\dd[zJ_1(z)]/\dd z=zJ_0(z)$。令$z=qr$、$r\dd r=z\dd z/q^2$，上下限由$0,a$变为$0,qa$，得到
\[
 U=\frac{2\pi C}{q^2}[zJ_1(z)]_0^{qa}
 =\frac{2\pi Ca}{q}J_1(qa).
\]
\Diagram{fig:o3circlea}{圆孔归一化远场强度；首零由贝塞尔函数确定}{
\draw[axis](0,0)--(5.1,0)node[right]{$qa$};
\draw[axis](0,0)--(0,2.9)node[above]{$I/I(0)$};
\draw[curve]plot coordinates{(0.000000,2.4500000) (0.015000,2.4494488) (0.030000,2.4477958) (0.045000,2.4450429) (0.060000,2.4411932) (0.075000,2.4362510) (0.090000,2.4302219) (0.105000,2.4231125) (0.120000,2.4149310) (0.135000,2.4056864) (0.150000,2.3953891) (0.165000,2.3840506) (0.180000,2.3716836) (0.195000,2.3583018) (0.210000,2.3439200) (0.225000,2.3285544) (0.240000,2.3122218) (0.255000,2.2949402) (0.270000,2.2767288) (0.285000,2.2576076) (0.300000,2.2375974) (0.315000,2.2167202) (0.330000,2.1949986) (0.345000,2.1724563) (0.360000,2.1491175) (0.375000,2.1250075) (0.390000,2.1001521) (0.405000,2.0745777) (0.420000,2.0483117) (0.435000,2.0213819) (0.450000,1.9938165) (0.465000,1.9656444) (0.480000,1.9368951) (0.495000,1.9075983) (0.510000,1.8777842) (0.525000,1.8474834) (0.540000,1.8167266) (0.555000,1.7855450) (0.570000,1.7539699) (0.585000,1.7220326) (0.600000,1.6897649) (0.615000,1.6571983) (0.630000,1.6243646) (0.645000,1.5912954) (0.660000,1.5580223) (0.675000,1.5245769) (0.690000,1.4909905) (0.705000,1.4572943) (0.720000,1.4235191) (0.735000,1.3896958) (0.750000,1.3558546) (0.765000,1.3220255) (0.780000,1.2882382) (0.795000,1.2545218) (0.810000,1.2209050) (0.825000,1.1874160) (0.840000,1.1540825) (0.855000,1.1209315) (0.870000,1.0879896) (0.885000,1.0552825) (0.900000,1.0228356) (0.915000,0.9906731) (0.930000,0.9588190) (0.945000,0.9272962) (0.960000,0.8961270) (0.975000,0.8653328) (0.990000,0.8349343) (1.005000,0.8049512) (1.020000,0.7754025) (1.035000,0.7463063) (1.050000,0.7176799) (1.065000,0.6895395) (1.080000,0.6619004) (1.095000,0.6347773) (1.110000,0.6081837) (1.125000,0.5821322) (1.140000,0.5566346) (1.155000,0.5317015) (1.170000,0.5073428) (1.185000,0.4835675) (1.200000,0.4603834) (1.215000,0.4377977) (1.230000,0.4158163) (1.245000,0.3944445) (1.260000,0.3736865) (1.275000,0.3535458) (1.290000,0.3340247) (1.305000,0.3151248) (1.320000,0.2968469) (1.335000,0.2791909) (1.350000,0.2621558) (1.365000,0.2457397) (1.380000,0.2299401) (1.395000,0.2147537) (1.410000,0.2001762) (1.425000,0.1862030) (1.440000,0.1728282) (1.455000,0.1600458) (1.470000,0.1478488) (1.485000,0.1362295) (1.500000,0.1251797) (1.515000,0.1146908) (1.530000,0.1047533) (1.545000,0.0953573) (1.560000,0.0864925) (1.575000,0.0781479) (1.590000,0.0703123) (1.605000,0.0629739) (1.620000,0.0561205) (1.635000,0.0497397) (1.650000,0.0438187) (1.665000,0.0383442) (1.680000,0.0333029) (1.695000,0.0286811) (1.710000,0.0244650) (1.725000,0.0206406) (1.740000,0.0171935) (1.755000,0.0141095) (1.770000,0.0113742) (1.785000,0.0089731) (1.800000,0.0068915) (1.815000,0.0051150) (1.830000,0.0036291) (1.845000,0.0024191) (1.860000,0.0014707) (1.875000,0.0007695) (1.890000,0.0003013) (1.905000,0.0000519) (1.920000,0.0000074) (1.935000,0.0001540) (1.950000,0.0004782) (1.965000,0.0009666) (1.980000,0.0016061) (1.995000,0.0023838) (2.010000,0.0032871) (2.025000,0.0043038) (2.040000,0.0054219) (2.055000,0.0066296) (2.070000,0.0079157) (2.085000,0.0092692) (2.100000,0.0106794) (2.115000,0.0121362) (2.130000,0.0136295) (2.145000,0.0151499) (2.160000,0.0166883) (2.175000,0.0182360) (2.190000,0.0197847) (2.205000,0.0213265) (2.220000,0.0228539) (2.235000,0.0243598) (2.250000,0.0258376) (2.265000,0.0272810) (2.280000,0.0286843) (2.295000,0.0300419) (2.310000,0.0313489) (2.325000,0.0326008) (2.340000,0.0337932) (2.355000,0.0349225) (2.370000,0.0359853) (2.385000,0.0369784) (2.400000,0.0378994) (2.415000,0.0387459) (2.430000,0.0395160) (2.445000,0.0402083) (2.460000,0.0408215) (2.475000,0.0413548) (2.490000,0.0418077) (2.505000,0.0421799) (2.520000,0.0424715) (2.535000,0.0426830) (2.550000,0.0428150) (2.565000,0.0428684) (2.580000,0.0428444) (2.595000,0.0427444) (2.610000,0.0425702) (2.625000,0.0423234) (2.640000,0.0420063) (2.655000,0.0416211) (2.670000,0.0411701) (2.685000,0.0406560) (2.700000,0.0400815) (2.715000,0.0394495) (2.730000,0.0387630) (2.745000,0.0380250) (2.760000,0.0372387) (2.775000,0.0364073) (2.790000,0.0355343) (2.805000,0.0346229) (2.820000,0.0336766) (2.835000,0.0326988) (2.850000,0.0316929) (2.865000,0.0306625) (2.880000,0.0296109) (2.895000,0.0285416) (2.910000,0.0274579) (2.925000,0.0263634) (2.940000,0.0252611) (2.955000,0.0241545) (2.970000,0.0230467) (2.985000,0.0219407) (3.000000,0.0208397) (3.015000,0.0197465) (3.030000,0.0186640) (3.045000,0.0175949) (3.060000,0.0165419) (3.075000,0.0155073) (3.090000,0.0144937) (3.105000,0.0135032) (3.120000,0.0125381) (3.135000,0.0116002) (3.150000,0.0106914) (3.165000,0.0098135) (3.180000,0.0089680) (3.195000,0.0081564) (3.210000,0.0073800) (3.225000,0.0066398) (3.240000,0.0059370) (3.255000,0.0052724) (3.270000,0.0046467) (3.285000,0.0040605) (3.300000,0.0035142) (3.315000,0.0030081) (3.330000,0.0025424) (3.345000,0.0021172) (3.360000,0.0017324) (3.375000,0.0013877) (3.390000,0.0010828) (3.405000,0.0008173) (3.420000,0.0005906) (3.435000,0.0004020) (3.450000,0.0002507) (3.465000,0.0001360) (3.480000,0.0000567) (3.495000,0.0000119) (3.510000,0.0000003) (3.525000,0.0000209) (3.540000,0.0000723) (3.555000,0.0001531) (3.570000,0.0002620) (3.585000,0.0003974) (3.600000,0.0005580) (3.615000,0.0007420) (3.630000,0.0009479) (3.645000,0.0011742) (3.660000,0.0014191) (3.675000,0.0016810) (3.690000,0.0019582) (3.705000,0.0022491) (3.720000,0.0025519) (3.735000,0.0028651) (3.750000,0.0031869) (3.765000,0.0035157) (3.780000,0.0038498) (3.795000,0.0041877) (3.810000,0.0045278) (3.825000,0.0048685) (3.840000,0.0052084) (3.855000,0.0055459) (3.870000,0.0058798) (3.885000,0.0062085) (3.900000,0.0065309) (3.915000,0.0068455) (3.930000,0.0071514) (3.945000,0.0074473) (3.960000,0.0077321) (3.975000,0.0080048) (3.990000,0.0082646) (4.005000,0.0085106) (4.020000,0.0087419) (4.035000,0.0089578) (4.050000,0.0091577) (4.065000,0.0093410) (4.080000,0.0095071) (4.095000,0.0096558) (4.110000,0.0097865) (4.125000,0.0098991) (4.140000,0.0099933) (4.155000,0.0100689) (4.170000,0.0101259) (4.185000,0.0101642) (4.200000,0.0101841) (4.215000,0.0101854) (4.230000,0.0101686) (4.245000,0.0101337) (4.260000,0.0100812) (4.275000,0.0100114) (4.290000,0.0099246) (4.305000,0.0098215) (4.320000,0.0097025) (4.335000,0.0095682) (4.350000,0.0094191) (4.365000,0.0092561) (4.380000,0.0090796) (4.395000,0.0088906) (4.410000,0.0086897) (4.425000,0.0084777) (4.440000,0.0082555) (4.455000,0.0080238) (4.470000,0.0077836) (4.485000,0.0075356) (4.500000,0.0072808)};
\node[left]at(0,2.45){$1$};
\draw[guide](1.915853,0)--(1.915853,.55);
\node[below]at(1.915853,0){$3.8317$};
\node[right]at(2.15,.45){首暗环};
}
轴上直接积分为$U(0)=C\pi a^2$，因此图\ref{fig:o3circlea}的归一化场和强度为
\[
 \frac{U}{U(0)}=\frac{2J_1(qa)}{qa},\qquad
 \frac I{I(0)}=\left[\frac{2J_1(qa)}{qa}\right]^2.
\]
$J_1(z)=z/2+O(z^3)$使轴上极限为1，不存在奇点。首个非零根$qa=3.8317$给
\[
 \sin\theta_0=\frac{3.8317}{2\pi}\frac\lambda a
 =1.2197\frac\lambda D,\quad D=2a.
\]
在小角度下这也是瑞利角间隔。两个互不相干点源的像应把两套强度相加；若相干，交叉项依相位而变，不能直接引用同一个双点强度判据。系数1.22由均匀圆孔的贝塞尔首零决定，矩形孔和高斯孔径需重新计算。
\end{solution}

\begin{example}\label{ex:o3gaussian}
孔面场振幅为$U(\xi)=U_0\ee^{-\xi^2/w^2}$，孔径足够大不截断高斯尾部。求远场强度及其$1/\ee^2$角半宽，与均匀硬边孔径比较。
\Diagram{fig:o3gaussq}{高斯振幅没有硬边界}{\draw[axis](-2.5,0)--(2.8,0)node[right]{$\xi/w$};\draw[axis](0,0)--(0,2.7)node[above]{$U/U_0$};\draw[curve,domain=-2.3:2.3,samples=80]plot(\x,{2.2*exp(-\x*\x)});\node[left]at(0,2.2){$1$};\draw(1,0)--(1,-.08)node[below]{$1$};}
\end{example}
\begin{solution}
图\ref{fig:o3gaussq}的远场积分为$F(q)=U_0\int_{-\infty}^{\infty}\ee^{-\xi^2/w^2-iq\xi}\dd\xi$。它正是第一讲已用分部积分及初值方程证明的式\eqref{eq:o1gaussintegral}。代入该式的$a=w^{-2}>0$、$b=-iq$，有$\sqrt{\pi/a}=w\sqrt\pi$、$b^2/(4a)=-q^2w^2/4$，因此
\[
 F(q)=U_0w\sqrt\pi\ee^{-q^2w^2/4}.
\]
轴上$q=0$时恢复孔面场的积分$U_0w\sqrt\pi$；方向相关的衰减完全由指数中的$q^2w^2$决定。
\Diagram{fig:o3gaussa}{高斯远场无旁瓣，角宽与孔面场宽成反比}{\draw[axis](-2.6,0)--(2.8,0)node[right]{$\theta/\theta_d$};\draw[axis](0,0)--(0,2.7)node[above]{$I/I(0)$};\draw[curve,domain=-2.4:2.4,samples=80]plot(\x,{2.3*exp(-2*\x*\x)});\draw[guide](-2.4,.3113)--(2.4,.3113);\node[left]at(0,.3113){$\ee^{-2}$};\node[left]at(0,2.3){$1$};\draw[guide](1,0)--(1,.3113);\draw(1,0)--(1,-.08)node[below]{$1$};}
图\ref{fig:o3gaussa}中的强度比为$\exp(-q^2w^2/2)$。取$I/I(0)=\ee^{-2}$，则$|q|=2/w$；近轴$q\simeq k\theta$，所以$1/\ee^2$角半宽为
\[
 \theta_d=\frac2{kw}=\frac\lambda{\pi w}.
\]
题中的$w$是场下降至$1/\ee$、强度下降至$1/\ee^2$的孔面半宽。若给的是强度高斯式$\exp(-\xi^2/w_I^2)$，须先取场的平方根，不能沿用相同的$w$。平滑高斯孔面对应无有限零点、无旁瓣的高斯远场，宽度须按约定强度比定义。$1/\ee^2$半宽与硬孔首零半宽采用不同的宽度标准。
\end{solution}

\begin{example}\label{ex:o3twoDgrating}
二维矩形周期阵列沿$x,y$周期为$b_x,b_y$，分别有$N_x,N_y$个相同小孔，正常入射。求主极大的方向条件及两方向首零宽度。
\Diagram{fig:o3twoDq}{二维有限孔阵列，图只显示部分周期}{\foreach\x in{0,1,2,3}{\foreach\y in{0,1,2}{\draw[curve](\x,\y)circle(.13);}}\draw[axis](-.5,-.5)--(4,-.5)node[right]{$x$};\draw[axis](-.5,-.5)--(-.5,3)node[above]{$y$};}
\end{example}
\begin{solution}
令$k_0=2\pi/\lambda$。图\ref{fig:o3twoDq}的平移相位可分离为
\[
 \sum_{j=0}^{N_x-1}\ee^{-ik_0b_xl_xj}
 \sum_{\ell=0}^{N_y-1}\ee^{-ik_0b_yl_y\ell}.
\]
两方向同时相长要求$b_xl_x=m\lambda$、$b_yl_y=n\lambda$。
\Diagram{fig:o3twoDa}{允许主峰位于方向余弦平面内的二维格点}{\draw[axis](-2,0)--(2.4,0)node[right]{$l_x$};\draw[axis](0,-1.8)--(0,2.2)node[above]{$l_y$};\draw[guide](0,0)circle(1.65);\foreach\x in{-1,0,1}{\foreach\y in{-1,0,1}{\fill(\x,\y)circle(1.5pt);}}}
还须$l_x^2+l_y^2\leq1$，图\ref{fig:o3twoDa}圆外格点不对应传播波。固定另一方向，首零宽度为$\Delta l_x=\lambda/(N_xb_x)$、$\Delta l_y=\lambda/(N_yb_y)$。单孔包络再乘到阵列因子上，某格点可能因孔形而缺级。二维条件不能只用总孔数$N_xN_y$替代一个方向的有效周期数。
\end{solution}

\begin{example}\label{ex:o3unequalresolution}
两条互不相干谱线的单线分布为偶函数$F$，强度比为$1:r$，峰中心为$\pm a$。解释为何等强时中点曲率判据不能直接用于$r\ne1$，写出真正峰合并边界的方程。
\Diagram{fig:o3unequalq}{弱谱线可能被强谱线的旁瓣或翼部覆盖}{\draw[axis](-3,0)--(3.3,0)node[right]{$x$};\draw[curve,domain=-2.8:2.8,samples=90]plot(\x,{2.3*exp(-2*(\x+1)^2)});\draw[guide,domain=-2.8:2.8,samples=90]plot(\x,{.7*exp(-2*(\x-1)^2)});}
\end{example}
\begin{solution}
图\ref{fig:o3unequalq}总强度$G(x)=F(x+a)+rF(x-a)$。当$r\ne1$时$G$不再关于0对称，$G'(0)=(1-r)F'(a)$一般不为零。因此即使$G''(0)=0$，中点也不是极值，不能用它判定两峰是否分开。
\Diagram{fig:o3unequala}{不等强高斯线形在临界间距处的合强度，$r=0.3$}{\draw[axis](-3.6,0)--(3.7,0)node[right]{$x$};\draw[axis](0,0)--(0,2.9)node[above]{$G$};\draw[curve]plot coordinates {(-3.40000,0.008830) (-3.35976,0.010673) (-3.31953,0.012859) (-3.27929,0.015443) (-3.23905,0.018487) (-3.19882,0.022059) (-3.15858,0.026235) (-3.11834,0.031102) (-3.07811,0.036752) (-3.03787,0.043289) (-2.99763,0.050823) (-2.95740,0.059476) (-2.91716,0.069376) (-2.87692,0.080663) (-2.83669,0.093484) (-2.79645,0.107991) (-2.75621,0.124347) (-2.71598,0.142717) (-2.67574,0.163271) (-2.63550,0.186182) (-2.59527,0.211622) (-2.55503,0.239760) (-2.51479,0.270761) (-2.47456,0.304783) (-2.43432,0.341970) (-2.39408,0.382454) (-2.35385,0.426349) (-2.31361,0.473745) (-2.27337,0.524709) (-2.23314,0.579277) (-2.19290,0.637453) (-2.15266,0.699204) (-2.11243,0.764459) (-2.07219,0.833103) (-2.03195,0.904978) (-1.99172,0.979877) (-1.95148,1.057549) (-1.91124,1.137690) (-1.87101,1.219952) (-1.83077,1.303937) (-1.79053,1.389206) (-1.75030,1.475273) (-1.71006,1.561619) (-1.66982,1.647686) (-1.62959,1.732891) (-1.58935,1.816630) (-1.54911,1.898282) (-1.50888,1.977220) (-1.46864,2.052818) (-1.42840,2.124458) (-1.38817,2.191542) (-1.34793,2.253499) (-1.30769,2.309791) (-1.26746,2.359927) (-1.22722,2.403466) (-1.18698,2.440025) (-1.14675,2.469287) (-1.10651,2.491006) (-1.06627,2.505012) (-1.02604,2.511211) (-0.98580,2.509591) (-0.94556,2.500220) (-0.90533,2.483249) (-0.86509,2.458905) (-0.82485,2.427492) (-0.78462,2.389386) (-0.74438,2.345025) (-0.70414,2.294912) (-0.66391,2.239598) (-0.62367,2.179680) (-0.58343,2.115789) (-0.54320,2.048583) (-0.50296,1.978736) (-0.46272,1.906932) (-0.42249,1.833851) (-0.38225,1.760163) (-0.34201,1.686518) (-0.30178,1.613541) (-0.26154,1.541821) (-0.22130,1.471903) (-0.18107,1.404289) (-0.14083,1.339426) (-0.10059,1.277707) (-0.06036,1.219464) (-0.02012,1.164969) (0.02012,1.114433) (0.06036,1.068004) (0.10059,1.025767) (0.14083,0.987750) (0.18107,0.953924) (0.22130,0.924203) (0.26154,0.898455) (0.30178,0.876497) (0.34201,0.858111) (0.38225,0.843037) (0.42249,0.830988) (0.46272,0.821651) (0.50296,0.814695) (0.54320,0.809771) (0.58343,0.806528) (0.62367,0.804608) (0.66391,0.803657) (0.70414,0.803329) (0.74438,0.803290) (0.78462,0.803223) (0.82485,0.802831) (0.86509,0.801841) (0.90533,0.800006) (0.94556,0.797108) (0.98580,0.792959) (1.02604,0.787405) (1.06627,0.780321) (1.10651,0.771619) (1.14675,0.761239) (1.18698,0.749154) (1.22722,0.735369) (1.26746,0.719914) (1.30769,0.702847) (1.34793,0.684251) (1.38817,0.664228) (1.42840,0.642900) (1.46864,0.620404) (1.50888,0.596890) (1.54911,0.572517) (1.58935,0.547451) (1.62959,0.521859) (1.66982,0.495912) (1.71006,0.469776) (1.75030,0.443616) (1.79053,0.417588) (1.83077,0.391839) (1.87101,0.366507) (1.91124,0.341720) (1.95148,0.317590) (1.99172,0.294219) (2.03195,0.271694) (2.07219,0.250087) (2.11243,0.229459) (2.15266,0.209856) (2.19290,0.191309) (2.23314,0.173839) (2.27337,0.157456) (2.31361,0.142157) (2.35385,0.127930) (2.39408,0.114756) (2.43432,0.102606) (2.47456,0.091446) (2.51479,0.081237) (2.55503,0.071934) (2.59527,0.063491) (2.63550,0.055858) (2.67574,0.048984) (2.71598,0.042817) (2.75621,0.037306) (2.79645,0.032398) (2.83669,0.028046) (2.87692,0.024200) (2.91716,0.020813) (2.95740,0.017843) (2.99763,0.015247) (3.03787,0.012987) (3.07811,0.011026) (3.11834,0.009331) (3.15858,0.007871) (3.19882,0.006618) (3.23905,0.005546) (3.27929,0.004633) (3.31953,0.003858) (3.35976,0.003202) (3.40000,0.002649)};\draw[guide](0.14049,0.803290)--(1.54049,0.803290);\fill(0.74049,0.803290)circle(1.5pt)node[above right]{$x_*$};}
应联立$F'(x_*+a)+rF'(x_*-a)=0$和$F''(x_*+a)+rF''(x_*-a)=0$，解未知$x_*,a$。例如$F(x)=\ee^{-x^2}$，两次求导为$F'=-2x\ee^{-x^2}$、$F''=(-2+4x^2)\ee^{-x^2}$。记$u=x_*+a,v=x_*-a$，第一式给$r\ee^{-v^2}=-(u/v)\ee^{-u^2}$；代入第二式并乘$v$，得到
\[
 v(4u^2-2)-u(4v^2-2)=2(u-v)(2uv+1)=0.
\]
非零分离$a>0$对应$uv=-1/2$，即$a^2-x_*^2=1/2$。再代回第一式，
\[
 r=\frac{a+x_*}{a-x_*}\ee^{-4ax_*}.
\]
对$r=0.3$取弱峰一侧$x_*>0$，数值求这两个标量方程，得$a=1.023881$、$x_*=0.740495$。图\ref{fig:o3unequala}按该解画出峰、谷合并处的水平切线。等强$r=1$时对称性给$x_*=0$，才退回中点曲率条件。
\end{solution}
